\section*{الفصل \ref{ch:g10:reaction-progress-table} --- جدول تقدّم التفاعل}

\begin{solution}{exo:g10:reaction-progress-table:1}
\ce{H2}: $4.0 - 2x$؛ \ce{O2}: $3.0 - x$؛ \ce{H2O}: $2x$. ينفد ثنائي
الهيدروجين عند $x = 2.0$، وثنائي الأكسجين عند $x = 3.0$: $x_{\max} = \qty{2.0}{mol}$،
وثنائي الهيدروجين هو المحدِّد. \omterm{def:g10:reaction-progress-table:system}{الحالة النهائية}: \qty{0}{mol} \ce{H2}،
و\qty{1.0}{mol} \ce{O2}، و\qty{4.0}{mol} \ce{H2O}.
\end{solution}

\begin{solution}{exo:g10:reaction-progress-table:2}
ينفد الكربون عند $x = 3.0$، وثنائي الأكسجين عند $x = 2.0$:
$x_{\max} = \qty{2.0}{mol}$، وثنائي الأكسجين هو المحدِّد. \omterm{def:g10:reaction-progress-table:system}{الحالة النهائية}:
\qty{1.0}{mol} من الكربون، ولا ثنائي أكسجين، و\qty{2.0}{mol} من ثنائي أكسيد الكربون.
\end{solution}

\begin{solution}{exo:g10:reaction-progress-table:3}
$x_{\max} = \qty{0.30}{mol}$ (الكبريت محدِّد). \omterm{def:g10:reaction-progress-table:system}{الحالة النهائية}:
\qty{0.20}{mol} من الحديد، ولا كبريت، و\qty{0.30}{mol} من كبريتيد الحديد.
\end{solution}

\begin{solution}{exo:g10:reaction-progress-table:4}
ينفد ثنائي النيتروجين عند $x = 1.0$، وثنائي الهيدروجين عند $x = 2.4/3 = 0.80$:
$x_{\max} = \qty{0.80}{mol}$. \omterm{def:g10:reaction-progress-table:system}{الحالة النهائية}: \qty{0.20}{mol} \ce{N2}، ولا
\ce{H2}، و$2 \times 0.80 = \qty{1.6}{mol}$ \ce{NH3}.
\end{solution}

\begin{solution}{exo:g10:reaction-progress-table:5}
(ب): $4/2 = 2/1$. وفي (أ) أول أكسيد الكربون محدِّد، وكذلك في (ج).
\end{solution}

\begin{solution}{exo:g10:reaction-progress-table:6}
$n(\ce{CH4}) = 32 / 16.0 = \qty{2.0}{mol}$، إذن $x_{\max} = \qty{2.0}{mol}$
\omterm{def:g10:reaction-progress-table:stoichiometric}{لخليط ستوكيومتري}. ثنائي الأكسجين: $2 \times 2.0 = \qty{4.0}{mol}$،
أي $4.0 \times 32.0 = \qty{128}{g}$. الماء:
$\qty{4.0}{mol} \times \qty{18.0}{g/mol} = \qty{72}{g}$.
\end{solution}

\begin{solution}{exo:g10:reaction-progress-table:7}
$n(\ce{C3H8}) = 1000 / 44.0 = \qty{22.7}{mol}$؛ ثنائي أكسيد الكربون
$3 \times 22.7 = \qty{68.2}{mol}$؛ $V = 68.2 \times 24.1 =
\qty{1.64e3}{L}$، أي نحو \qty{1.6}{m^3}.
\end{solution}

\begin{solution}{exo:g10:reaction-progress-table:8}
عند $x = \qty{1.0}{mol}$: \qty{1.0}{mol} \ce{CH4}، و\qty{1.0}{mol}
\ce{O2}، و\qty{1.0}{mol} \ce{CO2}، و\qty{2.0}{mol} \ce{H2O}. ويُستهلك ثنائي
الأكسجين عند $x = \qty{1.5}{mol}$: فيتوقف التفاعل هناك، ولا تُبلغ
التقدّمات الأكبر أبدًا.
\end{solution}

\begin{solution}{exo:g10:reaction-progress-table:9}
$n_0(\ce{Zn}) = 1.30 / 65.4 = \qty{0.0199}{mol}$؛
$n_0(\ce{Cu^{2+}}) = 0.10 \times 0.1000 = \qty{0.0100}{mol}$. \omterm{def:g9:ions:ion}{أيونات}
النحاس هي المحدِّدة: $x_{\max} = \qty{0.0100}{mol}$. النحاس المتكوّن:
$0.0100 \times 63.5 = \qty{0.635}{g}$؛ والزنك المتبقي:
$(0.0199 - 0.0100) \times 65.4 = \qty{0.65}{g}$. لا تبقى \omterm{def:g9:ions:ion}{أيونات} نحاس:
فقد زال اللون الأزرق.
\end{solution}

\begin{solution}{exo:g10:reaction-progress-table:10}
$n_0(\ce{Mg}) = 0.48 / 24.3 = \qty{0.0198}{mol}$، وكان سينفد عند
$x = 0.0198$؛ و$n_0(\ce{H+}) = \qty{0.020}{mol}$، وكانت ستنفد عند
$x = 0.010$. فالحمض هو المحدِّد: $x_{\max} = \qty{0.010}{mol}$، إذن
\qty{0.010}{mol} من ثنائي الهيدروجين، أي $0.010 \times 24.1 = \qty{0.24}{L}$.
\end{solution}

\begin{solution}{exo:g10:reaction-progress-table:11}
\omterm{def:g10:reaction-progress-table:stoichiometric}{الخليط الستوكيومتري} مع \qty{0.30}{mol} من الكبريت يحتاج إلى
\qty{0.30}{mol} من الحديد؛ وهنا \qty{0.20}{mol} زيادة، أي
$0.20 / 0.30 \approx \qty{67}{\percent}$ بزيادة.
\end{solution}

\begin{solution}{exo:g10:reaction-progress-table:12}
كميتان متساويتان: $m(\ce{Fe})/55.8 = m(\ce{S})/32.1$، مع
$m(\ce{Fe}) + m(\ce{S}) = \qty{10.0}{g}$. إذن
$m(\ce{Fe}) = 10.0 \times 55.8 / (55.8 + 32.1) = \qty{6.35}{g}$
و$m(\ce{S}) = \qty{3.65}{g}$.
\end{solution}

\begin{solution}{exo:g10:reaction-progress-table:13}
\begin{enumerate}
  \item $n_0(\ce{CaCO3}) = 10.0 / 100.1 = \qty{0.0999}{mol}$، وهو \omterm{def:g7:chemical-reactions:reactant}{المتفاعل}
        الوحيد: $x_{\max} = \qty{0.0999}{mol}$؛ فيعطي
        \qty{0.0999}{mol} من \ce{CaO} ومن \ce{CO2}، أي
        $0.0999 \times 24.1 = \qty{2.41}{L}$ من ثنائي أكسيد الكربون.
  \item $n_0(\ce{CaO}) = \qty{0.0999}{mol}$، $n_0(\ce{H2O}) = 1.8 / 18.0
        = \qty{0.100}{mol}$: فالجير الحي هو المحدِّد (بالكاد)؛ \omterm{def:g10:reaction-progress-table:stoichiometric}{والخليط
        ستوكيومتري} تقريبًا. \ce{Ca(OH)2}:
        $M = 40.1 + 2 \times (16.0 + 1.0) = \qty{74.1}{g/mol}$، والكتلة
        $0.0999 \times 74.1 = \qty{7.40}{g}$.
\end{enumerate}
\end{solution}

\begin{solution}{exo:g10:reaction-progress-table:14}
كمية \omterm{def:g9:ions:ion}{أيونات} الهيدروجين في كل دورق $n_0(\ce{H+}) = \qty{0.050}{mol}$؛ فهي تنفد عند
$x = \qty{0.025}{mol}$. المغنيسيوم: 0.15 و0.30 و0.45 و0.60 و0.75
و\qty{0.90}{g} تعطي 0.0062 و0.0123 و0.0185 و0.0247 و0.0309
و\qty{0.0370}{mol}. في الدوارق الأربعة الأولى يكون المغنيسيوم محدِّدًا
وثنائي الهيدروجين $n(\ce{Mg}) \times 24.1$: 0.15 و0.30 و0.45
و\qty{0.60}{L}. وفي الدورقين الأخيرين يكون الحمض محدِّدًا: $0.025 \times 24.1
= \qty{0.60}{L}$. والمنحنى خط مستقيم يمر بالمبدأ حتى
نحو \qty{0.61}{g} من المغنيسيوم، ثم استقرار أفقي عند
\qty{0.60}{L}.
\end{solution}

\begin{solution}{exo:g10:reaction-progress-table:15}
$n_0(\text{الإيثانول}) = 4.6 / 46.0 = \qty{0.100}{mol}$ (ينفد عند
$x = 0.100$)؛ و$n_0(\ce{O2}) = 4.8 / 24.1 = \qty{0.199}{mol}$ (ينفد
عند $x = 0.199/3 = 0.0664$). فثنائي الأكسجين هو المحدِّد،
$x_{\max} = \qty{0.0664}{mol}$. الإيثانول المتبقي: $0.100 - 0.0664 =
\qty{0.0336}{mol}$، أي $0.0336 \times 46.0 = \qty{1.5}{g}$. ثنائي أكسيد
الكربون: $2 \times 0.0664 = \qty{0.133}{mol}$، أي
$0.133 \times 24.1 = \qty{3.2}{L}$.
\end{solution}

\begin{solution}{pb:g10:reaction-progress-table:1}
\textbf{1.} اليسار: 2 Na، و6 N. اليمين: 2 Na، و$3 \times 2 = 6$ N. موزونة.

\textbf{2.} اليسار: 10 Na، و2 K، و2 N، و6 O. اليمين: K: 2؛ وNa: $5 \times 2 =
10$؛ وO: $1 + 5 = 6$؛ وN: 2. موزونة.

\textbf{3.} فصل \omterm{def:g9:periodic-table-first-look:periodic-table}{الجدول الدوري} (\cref{ch:g9:periodic-table-first-look}):
فالصوديوم يتفاعل بعنف مع الماء، فيعطي \omterm{def:g2:mixing-and-dissolving:dissolve}{محلولًا} أكّالًا وثنائي
الهيدروجين. وبعد الحادث تلامس الوسادة الجلد والعينين \omterm{def:g4:air-a-mixture-of-gases:air}{والهواء} الرطب
المزفور.

\textbf{4.} $M(\ce{NaN3}) = 23.0 + 3 \times 14.0 = \qty{65.0}{g/mol}$؛
$M(\ce{KNO3}) = 39.1 + 14.0 + 3 \times 16.0 = \qty{101.1}{g/mol}$.

\textbf{5.} \ce{NaN3}: $1.00 - 2x$؛ \ce{Na}: $2x$؛ \ce{N2}: $3x$.

\textbf{6.} $1.00 - 2x = 0$ عند $x_{\max} = \qty{0.500}{mol}$؛ وأزيد
الصوديوم، \omterm{def:g7:chemical-reactions:reactant}{المتفاعل} الوحيد، هو المحدِّد.

\textbf{7.} \ce{Na}: $2 \times 0.500 = \qty{1.00}{mol}$؛ \ce{N2}:
$3 \times 0.500 = \qty{1.50}{mol}$.

\textbf{8.} $1.50 \times 24.1 = \qty{36.2}{L}$.

\textbf{9.} \ce{Na}: $1.00 - 10x$؛ \ce{KNO3}: $n - 2x$؛ \ce{K2O}: $x$؛
\ce{Na2O}: $5x$؛ \ce{N2}: $x$.

\textbf{10.} $1.00/10 = n/2$، إذن $n = \qty{0.200}{mol}$، أي
$0.200 \times 101.1 = \qty{20.2}{g}$.

\textbf{11.} $x_{\max} = \qty{0.100}{mol}$: لا صوديوم، ولا نترات بوتاسيوم،
و\qty{0.100}{mol} \ce{K2O}، و\qty{0.500}{mol} \ce{Na2O}،
و\qty{0.100}{mol} \ce{N2}.

\textbf{12.} كان الصوديوم سينفد عند $x = 0.100$، ونترات البوتاسيوم عند
$x = 0.150/2 = 0.075$: فالنترات هي المحدِّدة، $x_{\max} =
\qty{0.075}{mol}$، ويبقى في الوسادة $1.00 - 10 \times 0.075 = \qty{0.25}{mol}$
من \omterm{def:g9:periodic-table-first-look:metal}{فلز} الصوديوم، جاهزًا للتفاعل مع الرطوبة.

\textbf{13.} $1.50 + 0.100 = \qty{1.60}{mol}$ من النيتروجين لكل \omterm{def:g10:the-mole:amount}{مول} من
أزيد الصوديوم.

\textbf{14.} $n = 60 / 24.1 = \qty{2.49}{mol}$.

\textbf{15.} $2.49 / 1.60 = \qty{1.56}{mol}$ من أزيد الصوديوم.

\textbf{16.} $1.56 \times 65.0 = \qty{101}{g}$.

\textbf{17.} $0.200 \times 1.56 = \qty{0.311}{mol}$، أي
$0.311 \times 101.1 = \qty{31.5}{g}$ من نترات البوتاسيوم.

\textbf{18.} $2.49 / 1.50 = \qty{1.66}{mol}$، أي
$1.66 \times 65.0 = \qty{108}{g}$: فالتفاعل الثاني يوفّر نحو
\qty{7}{g} من أزيد الصوديوم.

\textbf{19.} \omterm{def:g10:reaction-progress-table:limiting}{التفاعل التام} لا يتوقف إلا حين يُستهلك متفاعله المحدِّد؛ وهنا أزيد
الصوديوم هو \omterm{def:g7:chemical-reactions:reactant}{المتفاعل} الوحيد، ففي \omterm{def:g10:reaction-progress-table:system}{الحالة النهائية} تكون كميته
$1.00 - 2x_{\max} = 0$. فلا يبقى في الوسادة شيء سام.

\textbf{20.} نحو \qty{101}{g} من أزيد الصوديوم تملأ وسادة هوائية سعتها
\qty{60}{L}.
\end{solution}
